% !TEX program = xelatex
% ============================================================
%  Arabic Communication Plan Template — Nibras Center
%  مشروع بحثي: خطة الاتصال
%  Compile with: xelatex main.tex (twice)
% ============================================================

\documentclass[11pt,a4paper,landscape]{article}

\usepackage[a4paper,landscape,margin=1.5cm,top=2cm,bottom=2cm,headheight=12pt]{geometry}
\usepackage{array}
\usepackage{booktabs}
\usepackage{tabularx}
\usepackage{enumitem}
\usepackage{float}
\usepackage{titlesec}
\usepackage{fancyhdr}
\usepackage{setspace}
\usepackage[table]{xcolor}

\usepackage{bidi}
\usepackage[no-math]{fontspec}
\usepackage[quiet]{polyglossia}

\setmainfont[Scale=1]{Amiri-Regular.ttf}
\newfontfamily\arabicfont[Script=Arabic,Scale=0.95]{Amiri-Regular.ttf}
\newfontfamily\englishfont[Scale=0.95]{TeX Gyre Termes}

\setdefaultlanguage[calendar=gregorian,hijricorrection=1]{arabic}
\setotherlanguage[variant=british]{english}

\usepackage[breaklinks,hidelinks]{hyperref}
\hypersetup{pdftitle={خطة الاتصال},pdfauthor={الباحث الرئيسي}}

\pagestyle{fancy}
\fancyhf{}
\fancyhead[R]{\small\itshape خطة الاتصال}
\fancyhead[L]{\small اسم المشروع}
\fancyfoot[C]{\small\thepage}
\renewcommand{\headrulewidth}{0.3pt}

\newcolumntype{L}[1]{>{\raggedright\arraybackslash}p{#1}}
\newcolumntype{C}[1]{>{\centering\arraybackslash}p{#1}}

\definecolor{headerColor}{HTML}{1A3C6B}

\setstretch{1.3}

\begin{document}

\begin{center}
{\LARGE\bfseries خطة الاتصال (\LR{Communication Plan})}\\[0.3em]
{\large عنوان المشروع البحثي — \today}\\[1em]
\end{center}

% ============================================================
% 1. STAKEHOLDER COMMUNICATION MATRIX
% ============================================================
\section{مصفوفة الاتصال مع أصحاب المصلحة}
\label{sec:matrix}

\begin{table}[H]
\centering
\footnotesize
\renewcommand{\arraystretch}{1.5}
\begin{tabularx}{\textwidth}{L{3cm} L{3.5cm} C{2cm} C{2cm} L{3cm} L{3cm} C{2cm}}
\toprule
\textbf{أصحاب المصلحة} & \textbf{المعلومات المطلوبة} & \textbf{التكرار} & \textbf{الوسيلة} & \textbf{المسؤول عن الإرسال} & \textbf{الصيغة} & \textbf{اللغة} \\
\midrule
الباحث الرئيسي & تقدم المهام والمخاطر & أسبوعياً & اجتماع + بريد & مدير المشروع & تقرير مختصر & عربية \\
\midrule
الباحثون المشاركون & المهام والنتائج & أسبوعياً & اجتماع & الباحث الرئيسي & محضر اجتماع & عربية \\
\midrule
الجهة الممولة & التقدم العام والميزانية & شهرياً & تقرير رسمي & الباحث الرئيسي & تقرير رسمي & عربية/إنجليزية \\
\midrule
اللجنة العلمية & النتائج والمنهجية & فصلياً & عرض تقديمي & الباحث الرئيسي & عرض + تقرير & عربية \\
\midrule
المجتمع الأكاديمي & نتائج البحث & عند النشر & ورقة بحثية & الباحث الرئيسي & ورقة محكمة & إنجليزية \\
\midrule
فني المختبر & التعليمات والمهام & يومياً & بريد + مباشر & باحث مشارك & تعليمات & عربية \\
\bottomrule
\end{tabularx}
\end{table}

% ============================================================
% 2. MEETING SCHEDULE
% ============================================================
\section{جدول الاجتماعات}
\label{sec:meetings}

\begin{table}[H]
\centering
\renewcommand{\arraystretch}{1.5}
\begin{tabularx}{\textwidth}{L{4cm} C{2.5cm} C{2.5cm} L{3cm} X}
\toprule
\textbf{نوع الاجتماع} & \textbf{التكرار} & \textbf{المدة} & \textbf{المشاركون} & \textbf{الهدف} \\
\midrule
اجتماع الفريق الأسبوعي & أسبوعياً (الاثنين) & 30 دقيقة & PI + Co + Tech & مراجعة التقدم وحل العقبات \\
\midrule
اجتماع المراجعة الشهري & شهرياً & 60 دقيقة & الفريق + PM & مراجعة شاملة للمشروع \\
\midrule
تقرير الجهة الممولة & شهرياً & --- & PI + Fund & تقرير كتابي عن التقدم \\
\midrule
عرض اللجنة العلمية & فصلياً & 90 دقيقة & PI + اللجنة & عرض النتائج والمنهجية \\
\midrule
اجتماع إغلاق المرحلة & عند نهاية كل مرحلة & 120 دقيقة & الفريق كاملاً & مراجعة المخرجات والتخطيط \\
\bottomrule
\end{tabularx}
\end{table}

% ============================================================
% 3. REPORTING SCHEDULE
% ============================================================
\section{جدول التقارير}
\label{sec:reporting}

\begin{table}[H]
\centering
\renewcommand{\arraystretch}{1.5}
\begin{tabularx}{\textwidth}{L{4cm} C{2.5cm} C{2.5cm} L{4cm} X}
\toprule
\textbf{نوع التقرير} & \textbf{التكرار} & \textbf{تاريخ التسليم} & \textbf{المستلمون} & \textbf{المحتوى} \\
\midrule
تقرير الحالة الأسبوعي & أسبوعياً & كل اثنين & PI, الفريق & تقدم المهام والعقبات \\
\midrule
تقرير الحالة الشهري & شهرياً & أول الشهر & PI, Fund, PM & ملخص شامل + ميزانية \\
\midrule
تقرير المرحلة & نهاية كل مرحلة & --- & Fund, اللجنة & مخرجات المرحلة \\
\midrule
التقرير النهائي & نهاية المشروع & --- & Fund, اللجنة & تقرير شامل + نتائج \\
\bottomrule
\end{tabularx}
\end{table}

% ============================================================
% 4. COMMUNICATION TOOLS
% ============================================================
\section{أدوات الاتصال}
\label{sec:tools}

\begin{itemize}
    \item \textbf{البريد الإلكتروني:} للمراسلات الرسمية والتقارير.
    \item \textbf{اجتماعات الفيديو:} \LR{Zoom/Teams} للاجتماعات عن بُعد.
    \item \textbf{منصة إدارة المشروع:} \LR{Trello/Asana} لتتبع المهام.
    \item \textbf{التخزين السحابي:} \LR{Google Drive/OneDrive} لمشاركة الملفات.
    \item \textbf{المراسلة الفورية:} \LR{WhatsApp/Slack} للتواصل السريع.
\end{itemize}

% ============================================================
% 5. ESCALATION PROCEDURE
% ============================================================
\section{إجراءات التصعيد}
\label{sec:escalation}

\begin{enumerate}
    \item \textbf{المستوى الأول:} العقبات اليومية تُحل في اجتماع الفريق الأسبوعي.
    \item \textbf{المستوى الثاني:} القضايا غير المحلولة تُرفع للباحث الرئيسي خلال 48 ساعة.
    \item \textbf{المستوى الثالث:} القضايا الحرجة تُرفع للجهة الممولة خلال أسبوع.
    \item \textbf{المستوى الرابع:} التغييرات الجوهرية تتطلب موافقة رسمية مكتوبة.
\end{enumerate}

\end{document}
